%% kapa: exam=1 round=29 year=2018 period=1 form=A subject=real_estate subject_ko=부동산학원론 seq=30 number=110
%% source: Q-Net 감정평가사 기출문제 (cst00309) artlSeq=5202808 file=extracted/1008/5202808/2018년 제29회 감정평가사 1차시험 기출문제/2018년 제29회 감정평가사 1차시험 기출문제/1교시 A형.pdf
%% conversion: kordoc 4.24.51 -> tools/gen1.py (auto); review=proofread-sonnet; flags=none
\begin{question}{110}
공시지가기준법에 의한 토지의 감정평가시 개별요인 세항목의 비교내용이 다음의 표와 같을 때 개별요인 비교치(격차율)는? (단, 주어진 자료 이외의 내용은 없음)
\par
{\small\centering
\begin{tabular}{|c|c|c|c|}\hline
조 건 & 항목 & 세항목 & 비교내용\\\hline
\multirow{2}{*}{접근조건} & \multirow{2}{*}{교통의 편부} & 취락과의 접근성 & 대상토지가 10\% 우세\\\cline{3-4}
 & & 농로의 상태 & 대상토지가 5\% 열세\\\hline
\multirow{2}{*}{자연조건} & 일조 등 & 일조, 통풍 등 & 대상토지가 10\% 우세\\\cline{2-4}
 & 토양, 토질 & 토양ㆍ토질의 양부 & 대상토지가 5\% 열세\\\hline
\multirow{2}{*}{획지조건} & 면적, 경사 등 & 경사도 & 대상토지가 5\% 열세\\\cline{2-4}
 & 경작의 편부 & \makecell{형상에 의한\\장애정도} & 동일함\\\hline
행정적조건 & \makecell{행정상의 조장\\및 규제 정도} & 용도지역 & 동일함\\\hline
기타조건 & 기타 & 장래의 동향 & 대상토지가 10\% 열세\\\hline
\end{tabular}\par}
\begin{choices}
  \item 0.980
  \item 0.955
  \item 0.950
  \item 0.943
  \item 0.934
\end{choices}
\end{question}
\answer{4}
